\documentclass[11pt,a4paper]{article}

\usepackage[margin=2cm]{geometry}
\usepackage{enumitem}
\usepackage[hidelinks]{hyperref}
\usepackage[T1]{fontenc}
\usepackage[utf8]{inputenc}
\usepackage{titlesec}

\setlength{\parindent}{0pt}
\pagenumbering{gobble}

\titleformat{\section}{\large\bfseries}{}{0em}{}[\titlerule]

\begin{document}

\begin{center}
{\Large \textbf{MIGUEL ANGEL ANAY GOMEZ}}\\
Cloud \& MLOps Engineer | Azure \& AWS | DevOps | LLMOps | Full Stack\\
\vspace{4pt}
\href{mailto:contacto@miguel-anay.nom.pe}{contacto@miguel-anay.com.pe} |
\href{https://www.linkedin.com/in/miguel-anay}{LinkedIn} |
Perú
|
{936327402} |
\href{https:/portfolio.miguel-anay.nom.pe}{Portfolio} 
\end{center}

\vspace{6pt}

\section*{Perfil Profesional}

Ingeniero de Software con más de \textbf{10 años de experiencia} en desarrollo backend, cloud y arquitectura de sistemas.  
Actualmente enfocado en \textbf{Cloud Engineering, DevOps, MLOps y LLMOps}, con experiencia real desplegando plataformas productivas en \textbf{AWS y Azure}, usando contenedores, CI/CD, infraestructura como código y agentes de IA para automatización y operación de soluciones escalables y seguras.

\vspace{6pt}

\section*{Experiencia Profesional}

\textbf{CEO \& Full Stack Developer – Yiwu Import Corporation} \hfill \textit{Abr 2025 – Actual}
\begin{itemize}[leftmargin=*]
  \item Fundador y líder técnico de plataformas digitales orientadas a facturación electrónica, comercio mayorista y automatización con IA.
  \item Diseño y despliegue de \textbf{FactuFacil}, sistema de facturación electrónica Headless con módulos de Compras, Ventas y Contabilidad.
  \item Integración directa con \textbf{SUNAT}, agentes de IA y sistemas externos como YAPE-CHECK.
  \item Plataforma desplegada en \textbf{AWS usando Docker} con \textbf{pipelines CI/CD} automatizados.
  \item Uso de \textbf{IA como MCP (Model Context Protocol)} y agentes para consultas, búsqueda de productos y gestión operativa.
  \item Desarrollo de \textbf{Cajon.pe}, plataforma de compras conjuntas a importadores mayoristas.
\end{itemize}

\textbf{Tecnologías:} React.js, Next.js, Vue.js, Angular, Docker, AWS, CI/CD, IA Generativa

\vspace{12pt}

\textbf{Software Engineer III – Encora (Cliente: Profuturo)} \hfill \textit{Nov 2020 – Abr 2025}
\begin{itemize}[leftmargin=*]
  \item Desarrollo de plataformas empresariales de alta disponibilidad para miles de usuarios.
  \item \textbf{Agencia Virtual}: plataforma de autoservicio que optimizó trámites de afiliados y redujo la carga operativa.
  \item \textbf{Clave Web}: sistema seguro de recuperación de credenciales que redujo incidencias de soporte en más del 40\%.
  \item \textbf{Disfruta Profuturo}: desarrollo frontend de portal de beneficios y promociones.
  \item Participación en despliegues productivos, pipelines CI/CD y mejora continua.
\end{itemize}

\textbf{Tecnologías:} React.js, Next.js, Vue.js, Angular, Docker

\vspace{12pt}

\textbf{Analista de Procesos – Prosegur} \hfill \textit{Abr 2018 – Mar 2020}
\begin{itemize}[leftmargin=*]
  \item \textbf{Web Audit}: plataforma web para auditoría de abastecimiento de ATMs mediante análisis de video.
  \item \textbf{Web PER}: sistema de monitoreo de rutas con GPS, mapas interactivos y alertas automáticas.
  \item \textbf{App de Supervisiones}: aplicación Android para captura de evidencias en campo en tiempo real.
\end{itemize}

\textbf{Tecnologías:} React.js, PHP Laravel, MySQL, SQL Server, Android Studio

\vspace{12pt}

\textbf{Desarrollador Web \& Full Stack .NET – Medios Industriales} \hfill \textit{Abr 2014 – Mar 2018}
\begin{itemize}[leftmargin=*]
  \item \textbf{Control de Tiempos Operativos (CTO)}: sistema interno para control preciso de jornadas laborales.
  \item \textbf{Balizas}: plataforma de monitoreo GPS en tiempo real de unidades blindadas con alertas inteligentes.
  \item Optimización de procesos operativos y mejora de la productividad empresarial.
\end{itemize}

\vspace{6pt}

\section*{Cloud, DevOps y MLOps}

\begin{itemize}[leftmargin=*]
  \item Diseño y operación de soluciones en \textbf{AWS y Azure}.
  \item Despliegue de aplicaciones y servicios con \textbf{Docker}.
  \item Implementación de \textbf{CI/CD} con GitHub Actions y Azure DevOps.
  \item Infraestructura como Código: \textbf{Terraform y Bicep}.
  \item Gestión de identidades, secretos y seguridad (Key Vault, RBAC, MSI).
  \item Observabilidad, monitoreo y troubleshooting end-to-end.
  \item Experiencia práctica en \textbf{LLMOps}: versionado de prompts, agentes, endpoints y flujos dev/test/prod.
\end{itemize}

\vspace{6pt}

\section*{Tecnologías}

\textbf{Cloud:} AWS, Azure (App Services, Functions, Container Apps, Storage, Azure ML)\\
\textbf{CI/CD:} GitHub Actions, Azure DevOps\\
\textbf{Contenedores:} Docker\\
\textbf{Observabilidad:} Application Insights, Log Analytics\\
\textbf{Lenguajes:} C\#, .NET, Python, SQL, JavaScript\\
\textbf{Frontend:} React, Next.js, Vue, Angular\\
\textbf{Otros:} Git, Kubernetes (básico), GitOps (básico)

\vspace{6pt}
\section*{Educación \& Especialización}

\textbf{Universidad Nacional de Ingeniería (UNI)}\\
Ingeniero Industrial \hfill \textit{2004 – 2012}\\
CIP: 301084

\vspace{4pt}

\textbf{Universidad Privada del Norte}\\
Diplomado en Gerencia de TI con mención en Disrupción Tecnológica\\
\textit{Ene 2025 – Sep 2025}

\vspace{6pt}

\textbf{Especialización Técnica y Certificaciones}

\begin{itemize}[leftmargin=*]
  \item \textbf{Advanced Data Engineer} – DMC Instituto (Diplomado)\\
  \textit{Sep 2025 – Actual | 144 horas | Certificado en proceso}
  
  \item \textbf{Gen AI Training Path} – Coursera Technical Track\\
  \textit{Mar 2025 | 60 horas | Certificación Credly}
  
  \item \textbf{Ingeniero DevOps – Azure Fundamentals AZ-900}\\
  Udemy / Microsoft Azure Certification\\
  \textit{Nov 2024 | 9 horas}
  
  \item \textbf{Arquitectura Hexagonal C\#.NET – Technical Track}\\
  \textit{Oct 2025 | 8.5 horas}
  
  \item \textbf{Python Data Science} – DSRP Certification Program\\
  \textit{Oct 2025 | 20 horas}
  
  \item \textbf{Experto en Business Intelligence}\\
  Sistemas UNI – Módulo de Base de Datos\\
  \textit{Feb 2016 | 48 horas}
  
  \item \textbf{Administración de SQL Server}\\
  Sistemas UNI – Módulo de Base de Datos\\
  \textit{Jun 2015 | 96 horas}
  
  \item \textbf{Software Developer Full Stack}\\
  Especialización Técnica
  
  \item \textbf{Visual Basic .NET}\\
  Institución de Formación Técnica\\
  \textit{Sep 2011 | 48 horas}
\end{itemize}

\section*{Idiomas}

Español: Nativo\\
Inglés: Técnico (Intermedio)

\end{document}
